\chapter{The Cycle Class Map}\label{cycle-class}
\chapterstatus{complete}

\section{Introduction}\label{cycle-class:section-introduction}

A subvariety $Z$ of codimension $p$ in a complex projective manifold $X$ has a cohomology class $[Z] \in H^{2p}(X; \ZZ)$. Topologically it is the
Poincar\'e dual of the fundamental class of $Z$, which exists although $Z$ may be singular. Analytically it is represented by the current of integration
over $Z$. The resulting \emph{cycle class map} $\mathrm{cl} \colon \CH^p(X) \to H^{2p}(X; \ZZ)$ is a ring homomorphism compatible with pullbacks, pushforwards
and Chern classes. Its image consists of Hodge classes (Corollary~\ref{cycle-class:corollary-hodge-classes}), and the Hodge conjecture asks whether, after tensoring
with $\QQ$, every Hodge class is in the image. This chapter constructs the map, states its properties, and shows that the conjecture can equivalently be
phrased with Chern classes of vector bundles. References: \cite[Chapter 19]{fulton-intersection}, \cite[Chapter 11]{voisin-hodge-1}, \cite[Chapter 0, Section 4]{griffiths-harris}.

\noindent\textbf{Prerequisites.} Chapter~\ref{intersection-theory} (Algebraic Cycles and Chow Groups), Chapter~\ref{gaga} (Analytic Spaces and GAGA),
Chapter~\ref{de-rham} (de Rham Cohomology), Chapter~\ref{poincare-duality} (Topological Manifolds and Poincar\'e Duality) and Chapter~\ref{hodge-decomposition}
(The Hodge Decomposition).

Throughout, $X$ is a nonsingular complex variety of dimension $n$, identified with the complex manifold $X^{\mathrm{an}}$ (Chapter~\ref{gaga}), and cohomology is singular
cohomology of $X^{\mathrm{an}}$.

\section{Fundamental classes}\label{cycle-class:section-fundamental-classes}

\begin{lemma}\label{cycle-class:lemma-bm-vanishing}
Let $S$ be a closed complex analytic subset of complex dimension $\leq d$ of a complex manifold. Then $H^{BM}_i(S; \ZZ) = 0$ for $i > 2d$, and $H^{BM}_{2d}(S; \ZZ)$ is free on the
irreducible components of dimension $d$.
\end{lemma}

\begin{proof}[Sketch of proof]
Analytic sets admit triangulations compatible with a stratification by complex submanifolds, the strata of $S$ having real dimension $\leq 2d$
\cite[Chapter 19]{fulton-intersection}. So the locally finite chain complex of $S$ vanishes above degree $2d$, which gives the first statement. For the second, let
$\Sigma = \mathrm{Sing}(S) \cup S_{<d}$, where $S_{<d}$ is the union of the components of dimension $< d$. It has complex dimension $\leq d - 1$, and the long exact sequence
$H^{BM}_{2d}(\Sigma) \to H^{BM}_{2d}(S) \to H^{BM}_{2d}(S \setminus \Sigma) \to H^{BM}_{2d-1}(\Sigma)$ has outer terms zero by the first statement. The complex manifold $S \setminus \Sigma$ is
canonically oriented, of real dimension $2d$, with one connected component for each $d$-dimensional irreducible component of $S$ (the regular locus of an irreducible analytic set is
connected), and $H^{BM}_{2d}$ of an oriented $2d$-manifold is free on its connected components (Remark~\ref{poincare-duality:remark-borel-moore}).
\end{proof}

\begin{definition}\label{cycle-class:definition-cycle-class}
Let $Z \subseteq X$ be a subvariety of codimension $p$. Its \emph{fundamental class} $[Z] \in H^{BM}_{2n-2p}(Z; \ZZ)$ is the class corresponding to $Z$ in
Lemma~\ref{cycle-class:lemma-bm-vanishing}. Its \emph{cycle class} $\mathrm{cl}(Z) \in H^{2p}(X; \ZZ)$ is the image of $[Z]$ under
$H^{BM}_{2n-2p}(Z) \cong H^{2p}(X, X \setminus Z) \to H^{2p}(X)$ (Remark~\ref{poincare-duality:remark-fundamental-classes-subspaces}). It extends linearly to
$\mathrm{cl} \colon Z^p(X) \to H^{2p}(X; \ZZ)$. For $X$ compact, $\mathrm{cl}(Z)$ is the Poincar\'e dual of $i_*[Z] \in H_{2n-2p}(X; \ZZ)$.
\end{definition}

\begin{proposition}\label{cycle-class:proposition-current}
Let $X$ be compact. The image of $\mathrm{cl}(Z)$ in $H^{2p}(X; \RR) \cong H^{2p}_{\mathrm{dR}}(X)$ is the class of the current of integration over $Z_{\mathrm{reg}}$ of
Proposition~\ref{complex-manifolds:proposition-subvariety-class}: for every closed $(2n - 2p)$-form $\varphi$,
\[
\int_X\mathrm{cl}(Z) \wedge \varphi = \int_{Z_{\mathrm{reg}}}\varphi .
\]
\end{proposition}

\begin{proof}[Sketch of proof]
For $Z$ nonsingular this is Theorem~\ref{de-rham:theorem-currents}. In general both sides are determined by their restriction to the complement of $\Sigma = \mathrm{Sing}(Z)$. The
topological class is determined there by Lemma~\ref{cycle-class:lemma-bm-vanishing}, and the current by the fact that integration over $Z_{\mathrm{reg}}$ is a closed current whose
restriction to $X \setminus \Sigma$ is the integration current of a closed submanifold. Two closed currents of degree $2p$ that agree off an analytic set of real codimension
$> 2p$ in $X$ are equal (the difference would be a closed current supported on $\Sigma$ of degree below its codimension, which vanishes by the support theorem) \cite[Chapter 3]{griffiths-harris},
\cite[Section 11.1]{voisin-hodge-1}.
\end{proof}

\section{The cycle class map and its properties}\label{cycle-class:section-properties}

\begin{theorem}\label{cycle-class:theorem-cycle-class}
Let $X$ be a nonsingular complex variety.
\begin{enumerate}
\item $\mathrm{cl}$ factors through rational equivalence and gives homomorphisms $\mathrm{cl} \colon \CH^p(X) \to H^{2p}(X; \ZZ)$, which form a ring homomorphism
$\CH^\bullet(X) \to H^{2\bullet}(X; \ZZ)$ for the intersection product and the cup product.
\item It is compatible with pullback along morphisms of nonsingular varieties and with proper pushforward (Gysin maps in cohomology), and with exterior products.
\item $\mathrm{cl}(c_i(E) \cap [X]) = c_i(E)$, the topological Chern class of the underlying complex vector bundle.
\item For $X$ projective, $\deg\alpha = \int_X\mathrm{cl}(\alpha)$ for $\alpha \in \CH^n(X)$.
\end{enumerate}
\end{theorem}

This is \cite[Chapter 19]{fulton-intersection}. For (1), rational equivalence is generated by the fibres $\Gamma(0) - \Gamma(\infty)$ of a family
$\Gamma \subseteq \PP^1 \times X$ (Proposition~\ref{intersection-theory:proposition-equivalences}), and $\mathrm{cl}(\Gamma(0)) = \mathrm{cl}(\Gamma(\infty))$ because the two points of $\PP^1$
have the same class in $H^2(\PP^1; \ZZ)$ and $\mathrm{cl}(\Gamma(t)) = \mathrm{pr}_{2*}(\mathrm{cl}(\Gamma) \cup \mathrm{pr}_1^*\mathrm{cl}(t))$. Multiplicativity reduces, by the moving lemma or
deformation to the normal cone, to transverse intersections, where it is the statement that cup product is Poincar\'e dual to transverse intersection
(Remark~\ref{poincare-duality:remark-fundamental-classes-subspaces}).

\begin{example}\label{cycle-class:example-classes}
\begin{enumerate}
\item On $\PP^n$, $\mathrm{cl}(h^k)$ is the $k$-th power of the positive generator of $H^2(\CC\PP^n; \ZZ)$, and $\mathrm{cl} \colon \CH^\bullet(\PP^n) \to H^{2\bullet}(\CC\PP^n; \ZZ)$ is an isomorphism of
rings (Example~\ref{intersection-theory:example-chow-ring-projective} and Theorem~\ref{singular-cohomology:theorem-projective-spaces}).
\item For a divisor $D$, $\mathrm{cl}(D) = c_1(\mathcal{O}(D))$, by (3) applied to $E = \mathcal{O}(D)$ and $c_1(\mathcal{O}(D)) \cap [X] = [D]$. So on $\CH^1(X) = \Pic(X)$ the cycle class map is the first
Chern class of Theorem~\ref{holomorphic-bundles:theorem-exponential}.
\item For a point $x$ of a projective $X$, $\mathrm{cl}(x)$ is the generator of $H^{2n}(X; \ZZ) \cong \ZZ$ dual to the fundamental class. So $\mathrm{cl}$ on $\CH_0$ is the degree, and it loses all the
information in the huge kernel of Example~\ref{intersection-theory:example-chow-zero-surfaces}.
\end{enumerate}
\end{example}

\section{Cycle classes are Hodge classes}\label{cycle-class:section-hodge-classes}

\begin{definition}\label{cycle-class:definition-hodge-classes}
Let $X$ be a compact K\"ahler manifold. The group of \emph{rational Hodge classes} of degree $2p$ is $\mathrm{Hdg}^p(X) = H^{2p}(X; \QQ) \cap H^{p,p}(X)$, the intersection taken inside
$H^{2p}(X; \CC)$. The group of \emph{integral Hodge classes} $\mathrm{Hdg}^p(X; \ZZ)$ is the preimage of $H^{p,p}(X)$ under $H^{2p}(X; \ZZ) \to H^{2p}(X; \CC)$; it contains the torsion
subgroup.
\end{definition}

\begin{corollary}\label{cycle-class:corollary-hodge-classes}
For a nonsingular projective complex variety $X$, $\mathrm{cl}(\CH^p(X)) \subseteq \mathrm{Hdg}^p(X; \ZZ)$, and $\mathrm{cl}(\CH^p(X)) \otimes \QQ \subseteq \mathrm{Hdg}^p(X)$.
\end{corollary}

\begin{proof}
By Proposition~\ref{cycle-class:proposition-current}, the image of $\mathrm{cl}(Z)$ in $H^{2p}(X; \CC)$ is the class of the integration current over $Z$, which lies in $H^{p,p}(X)$ by
Proposition~\ref{hodge-decomposition:proposition-cycle-classes-type}. Extend linearly.
\end{proof}

\begin{remark}\label{cycle-class:remark-hodge-conjecture}
The \emph{Hodge conjecture} (Chapter~\ref{hodge-conjecture}) asserts that $\mathrm{cl}(\CH^p(X)) \otimes \QQ = \mathrm{Hdg}^p(X)$ for every nonsingular projective complex variety. The
integral version, $\mathrm{cl}(\CH^p(X)) = \mathrm{Hdg}^p(X; \ZZ)$, is false: Atiyah and Hirzebruch found torsion classes that are not algebraic \cite{atiyah-hirzebruch-1962}, and
Koll\'ar found non-torsion counterexamples (Chapter~\ref{failures}). For $p = 1$ the integral statement is true: this is the Lefschetz $(1,1)$-theorem
(Chapter~\ref{lefschetz-11}), which for projective $X$ follows from Lemma~\ref{kodaira:lemma-integral-11} together with Example~\ref{cycle-class:example-classes}(2) and GAGA.
\end{remark}

\section{Cycles and coherent sheaves}\label{cycle-class:section-chern-character}

\begin{proposition}\label{cycle-class:proposition-chern-character-cycle}
Let $X$ be a nonsingular projective variety and $Z \subseteq X$ a subvariety of codimension $p$. Then $\mathrm{ch}(\mathcal{O}_Z) = [Z] + (\text{terms in } \CH^{>p}(X) \otimes \QQ)$, where
$\mathrm{ch}$ is computed with a finite locally free resolution of $\mathcal{O}_Z$. Consequently the $\QQ$-subspace of $H^{2\bullet}(X; \QQ)$ spanned by the classes of algebraic cycles equals the
$\QQ$-subalgebra generated by the Chern classes of algebraic vector bundles on $X$.
\end{proposition}

\begin{proof}
The first statement is in \cite[Chapter 15]{fulton-intersection}: it follows from Grothendieck--Riemann--Roch (Theorem~\ref{intersection-theory:theorem-grr}) applied to a resolution
of singularities of $Z$, or directly from the fact that the lowest nonzero term of $\mathrm{ch}$ of a sheaf supported in codimension $p$ is its fundamental cycle. On a nonsingular projective
variety every coherent sheaf has a finite resolution by locally free sheaves \cite[Exercises III.6.8 and III.6.9]{hartshorne-ag}, so $\mathrm{ch}(\mathcal{O}_Z)$ is a polynomial with rational
coefficients in Chern classes of vector bundles. By descending induction on $p$, every $[Z]$ lies in the $\QQ$-algebra generated by Chern classes of vector bundles. Conversely, Chern
classes of vector bundles are classes of algebraic cycles, $c_i(E) \cap [X] \in \CH^i(X)$ (Theorem~\ref{cycle-class:theorem-cycle-class}(3)), and products of cycle classes are cycle
classes (Theorem~\ref{cycle-class:theorem-cycle-class}(1)).
\end{proof}

\begin{remark}\label{cycle-class:remark-kahler-version}
Proposition~\ref{cycle-class:proposition-chern-character-cycle} suggests versions of the Hodge conjecture for compact K\"ahler manifolds: with analytic cycles, with Chern classes of
holomorphic vector bundles, or with Chern classes of coherent analytic sheaves. All of these fail for suitable complex tori: Voisin constructed a
four-dimensional torus with a nonzero Hodge class in $H^4$ that is not a combination of Chern classes of coherent sheaves \cite{voisin-2002}. So projectivity is essential
(Chapter~\ref{failures}).
\end{remark}

\begin{remark}\label{cycle-class:remark-deligne-cohomology}
The cycle class map refines to the \emph{Deligne cycle class} with values in Deligne cohomology $H^{2p}_{\mathcal{D}}(X; \ZZ(p))$. Its restriction to homologically trivial cycles is the
Abel--Jacobi map to the intermediate Jacobian (Chapter~\ref{intermediate-jacobians}), which sees information about $\CH^p(X)$ invisible to cohomology, such as $\Pic^0$ for $p = 1$.
\end{remark}

\section{Exercises}\label{cycle-class:section-exercises}

\begin{exercise}\label{cycle-class:exercise-hypersurface-class}
Let $X \subseteq \PP^n$ be a hypersurface of degree $d$. Show that $\mathrm{cl}(X) = dh$ in $H^2(\CC\PP^n; \ZZ)$.
\end{exercise}

\begin{solution}
$[X] = \mathrm{div}(f/x_0^d) + d[H_0]$ in $Z^1(\PP^n)$ for a defining polynomial $f$, so $[X] = dh$ in $\CH^1(\PP^n)$, and $\mathrm{cl}(h)$ is the generator. Alternatively,
$\mathrm{cl}(X) = c_1(\mathcal{O}(X)) = c_1(\mathcal{O}(d)) = dh$ by Example~\ref{cycle-class:example-classes}(2).
\end{solution}

\begin{exercise}\label{cycle-class:exercise-diagonal}
Let $X$ be a nonsingular projective variety of dimension $n$ and $\Delta \subseteq X \times X$ the diagonal. Show that $\deg(\Delta \cdot \Delta) = \chi_{\mathrm{top}}(X)$, and more generally
the \emph{Lefschetz fixed point formula}: for a morphism $f \colon X \to X$ with graph $\Gamma_f$, $\deg(\Gamma_f \cdot \Delta) = \sum_i(-1)^i\operatorname{tr}(f^* \mid H^i(X; \QQ))$.
\end{exercise}

\begin{solution}
Choose a basis $(e_a)$ of $H^\bullet(X; \QQ)$ with dual basis $(e^a)$ for the Poincar\'e pairing, $\int_Xe_a \cup e^b = \delta_a^b$. By the K\"unneth formula,
$\mathrm{cl}(\Delta) = \sum_a(-1)^{\deg e_a}e_a \otimes e^a$ (the sign makes $\int_{X \times X}\mathrm{cl}(\Delta) \cup (\alpha \otimes \beta) = \int_X\alpha \cup \beta$). Then
$\mathrm{cl}(\Gamma_f) = (1 \times f)^*\mathrm{cl}(\Delta)$ up to the ordering of factors, and $\int\mathrm{cl}(\Gamma_f) \cup \mathrm{cl}(\Delta) = \sum_a(-1)^{\deg e_a}\int_Xf^*e_a \cup e^a$, which is the
alternating trace. For $f = \mathrm{id}$ the traces are the Betti numbers. Degrees of intersections are integrals of cup products by Theorem~\ref{cycle-class:theorem-cycle-class}(4).
\end{solution}

\begin{exercise}\label{cycle-class:exercise-torsion-hodge}
Show that every torsion class in $H^{2p}(X; \ZZ)$ is an integral Hodge class, and that for $X$ with $H^{2p}(X; \ZZ)$ torsion-free and $H^{2p}(X; \CC) = H^{p,p}(X)$ every integral class is a
Hodge class.
\end{exercise}

\begin{solution}
A torsion class maps to $0 \in H^{2p}(X; \CC)$, which lies in $H^{p,p}$. The second statement is immediate from the definition.
\end{solution}
